\subsection*{返回类型与类型推导}

返回值和输入一样，可以是带形状数据、标量，或者不返回结果（\co@void@）。

除了显式声明返回类型，也可以让编译器推断：

\begin{lstlisting}[language=croqtile,caption={返回类型与类型推导}]
__co__ auto foo() {
  f32 [7, 16] result;
  // ...
  return result;
}
\end{lstlisting}

这里，编译器根据 \co@result@ 将 \co@auto@ 推导为 \co@f32 [7, 16]@。当输出形状在函数体内推导时，这种写法可以避免修改推导式后还要同步修改返回声明。若固定返回形状是接口的一项要求，则适合显式声明类型。
